% ============================================================
%  Abstract
% ============================================================
\begin{abstract}
Curating web-scale image--text data usually amounts to thresholding a frozen
vision--language similarity, which presupposes that scores are comparable across
pairs. On the Attribute facet of a controlled benchmark, raw CLIP similarity orders
$31.9\%$ of all clean/noisy pair combinations backwards, although the same encoder
separates the two classes far above chance. We attribute part of this behaviour to
additive score components that are predictable from the image and text features and
largely unrelated to whether a pair matches. CRD (cross-modal residual debiasing),
a post-hoc calibration, fits two small additive heads on the frozen features and
subtracts a scaled estimate of the component from the standardised similarity; inference needs
one pair, no reference bank, and labels only to select one scalar. We also present
FGBench, $88{,}024$ reciprocal cycles over five factual facets from COCO and Visual
Genome, with silver labels. On the three reported Attribute backbones CRD raises
AUROC (CLIP ViT-B/16: $0.6806\to0.7187$; SigLIP2: $0.7026\to0.7509$); $21$ of $24$
paired comparisons survive Holm, or $17$ of $20$ independent comparisons once four
monotone GMM duplicates are excluded. Two parameter-free debiasing baselines do not
reproduce the Attribute gain. A logistic-regression probe trained on a
$32{,}000$-row labeled subset is $0.47$ and $1.60$ points higher on Attribute and
Quantity, and an MLP probe $7.62$ higher on the exploratory Action facet, at the cost
of about four times CRD's label budget; at CRD's budget the logistic-regression probe
falls $4.36$ points below CRD on Attribute. The method leaves cycle-level forced-choice
accuracy (2AFC) unchanged by construction and does not surpass CSLS on Quantity or
Entity.
\end{abstract}
